% ============================================================================ %
%
%           Šablona bakalářské/diplomové práce
%
% Autor:    Ing. Jozef Říha (4. květen 2006)
%           (některé komentáře převzaty z dokumentu Ivana Pomykacze)
%
% Verze:    2017-01-19, Ing. Pavel Tomášek (tomasek@fai.utb.cz)
%
% Kódování: UTF-8 (žluťoučký kůň úpěl ďábelšké ódy)
%
% Sazba:    pdflatex prace.tex && pdflatex prace.tex
%           (nutné dvakrát pro korektní vložení citací a jiných referencí),
%           v případě umístění literatury do externího bib souboru je třeba volat
%           pdflatex statement.tex && bibtex statement.tex && pdflatex statement.tex && pdflatex statement.tex
%
% Tip:      Ve správně vysázeném českém textu by na konci řádku neměla zůstant
%           samotná jednopísmenná předložka. Na takové místo se vkládá
%           nezalomitelná mezera pomocí symbolu ~. Existuje program, který umí
%           zpracovat celý TeX dokument najednou podle českých konvencí:
%           http://petr.olsak.net/ftp/olsak/vlna/
%
%           Pozor! Vzhledem k požadovanému standardu PDF/A nesmí obrázky obsahovat 
%           alfa kanál (průhlednost).
%
% ============================================================================ %


\documentclass[a4paper,12pt]{article}

% Definice vzhledu a nastavení se načítá z následujícího souboru (netřeba editovat)
\input{tex/UTB.tex}

% Uživatelské definice -- upravte dle požadavků
\nastavfakultu{FAI}
	% FAI  -- pro Fakultu aplikované informatiky
	% FAME -- pro Fakultu managementu a ekonomiky
	% FHS  -- pro Fakultu humanitních studií
	% FLKR -- pro Fakultu logistiky a krizového řízení
	% FMK  -- pro Fakutlu mutimediálních komunikací
	% FT   -- pro Fakultu technologickou
	% UNI  -- pro Univerzitní institut
\nastavtyp{BP}
	% BP   -- bakalářská práce
	% DP   -- diplomová práce
\nastavrok{2024}
	% zadejte rok místo "xxxx"
\nastavjazyk{CZ}
	% CZ   -- práce bude v českém jazyce
	% EN   -- práce bude v anglickém jazyce

% Lze přidat vertikalni odsazeni nad (prvni parametr) a pod (druhy parametr)
% obrázky, tabulky i rovnice/soustavy rovnic
\nastavmezerukolemobrazku{0mm}{0mm}
\nastavmezerukolemtabulek{0mm}{0mm}
\nastavmezerukolemrovnic{0mm}{0mm}

\nastavautora{Martin Tvaróg}
\nastavnazevcz{Technologie Blockchain pro ověřování dokumentů}
\nastavnazeven{Blockchain technology for document verification} % Jen u anglicky psané práce
\nastavabstraktcz{Tato bakalářská práce se věnuje zkoumání a aplikaci technologie Blockchain pro ověřování dokumentů. V první části práce je provedena literární rešerše se zaměřením na technologie blockchain a decentralizované aplikace, která nám pomáhá pochopit základní principy této problematiky. Následuje výběr nejvhodnějších technologií pro zpracování dat a vývoj aplikace. V práci je také stanoven postup pro ověřování, zpracování a ukládání dokumentů do blockchainu a představena metoda pro verifikaci dokumentů. Součástí práce je návrh uživatelsky přívětivého frontendu. Funkčnost navržené aplikace je pak demonstrována na zvoleném scénáři, přičemž zásadní pozornost je věnována testování a zabezpečení aplikace. Práce tak představuje komplexní pohled na možnosti využití blockchainu v oblasti ověřování dokumentů.}

\nastavabstrakten{This bachelor’s thesis focuses on investigating and applying Blockchain technology for document verification. The first part of the thesis includes a literature review, concentrating on blockchain technology and decentralized applications, which aids in understanding the fundamental principles of this issue. This is followed by the selection of the most appropriate technologies for data processing and application development. The thesis also establishes a procedure for verification, processing and storing documents into the blockchain and introduces a method for document verification. Part of the thesis is the design of a user-friendly frontend. The functionality of the proposed application is then demonstrated on a selected scenario, with significant attention paid to testing and security of the application. Thus, the thesis presents a comprehensive view of the possibilities of using blockchain in the field of document verification.}
\nastavklicovaslovacz{Blockchain, Ethereum, Ethereum Attestation Service, Decentralizované aplikace, Kryptografie}
\nastavklicovaslovaen{Blockchain, Ethereum, Ethereum Attestation Service, Decentralized Applications, Cryptography}

% Následující příkaz nastaví standard /A-1b
\aplikujpdfa


% ============================================================================ %
\begin{document}

\titulnistrana

\zadani

\prohlaseni

\abstraktaklicovaslova


% ============================================================================ %
\clearpage
\thispagestyle{empty}
Zde je místo pro případné poděkování, motto, úryvky knih, básní atp.


% ============================================================================ %
\obsah  % Obsah je generován automaticky


% ============================================================================ %
\OdsazovaniOdstavcuStart % Nastaví odsazování odstavců dle zvoleného jazyka

% ============================================================================ %
% Encoding: UTF-8 (žluťoučký kůň úpěl ďábelšké ódy)
% ============================================================================ %

% ============================================================================ %
\nn{Úvod}
Technologie blockchain, která získala globální pozornost díky svému použití v digitální měně Bitcoin, nabízí inovativní možnosti v různých oblastech, včetně ověřování dokumentů. Bitcoin představuje decentralizovaný systém pro transakce, který eliminuje potřebu třetí strany jako ověřovatele. Tento koncept decentralizace je klíčovým impulsem pro tuto práci, jejímž cílem je vytvořit decentralizovanou aplikaci pro ověřování dokumentů pomocí blockchainu a konkrétně Ethereum blockchainu.

Hlavním úkolem této práce je seznámení se s principem blockchainu, výběr a aplikace Ethereum blockchainu jako vhodné technologie pro zpracování dat a vývoj aplikace. Technologie blockchain je tu využita pro ověřování, zpracování a ukládání dokumentů do systému a na následné představení způsobu verifikace těchto dokumentů. Důležitá je i tvorba intuitivního uživatelského rozhraní tzv - frontednu, který umožnuje uživatelům aplikace s block chainem jednoduše iteragovat.

V teoretické části práce se budeme hlouběji zabývat funkcemi a principy blockchainu, roli kryptografie v blockchainu a také vlastnostmi a využitím Ethereum blockchainu a decentralizovaných aplikací. Budou představeny také nástroje jako Ethereum Attestation Service a MetaMask, které jsou klíčové a umožní nám realizovat naše cíle, a React, který bude použit pro vývoj frontendu aplikace.

Praktická část práce pokrývá návrh a implementaci samotné aplikace. Bude zde představen scénář, který ilustruje, jak může univerzita generovat atestace a k nim připojené dokumenty, a jak je může student předat zaměstnavateli k ověření pravosti pomocí naší aplikace. Tato aplikační demonstrace ukáže, jak může být blockchain efektivně využit k ověření pravosti dokumentů v decentralizovaném systému.

V závěru práce shrneme dosažené výsledky a zhodnotíme výkon aplikace. Pozornost bude věnována také diskusi o možných budoucích vylepšeních a dopadech našeho výzkumu. Cílem práce je tedy vytvořit transparentní, bezpečný a decentralizovaný systém pro ověřování dokumentů s využitím technologie blockchainu.



% ============================================================================ %
\cast{Teoretická část}

\section{Blockchain}
Blockchain je distribuovaná databáze, která udržuje neustále rostoucí seznam uspořádaných záznamů, nazývaných bloky. Tyto bloky jsou propojeny pomocí kryptografie. Každý blok obsahuje kryptografický otisk předchozího bloku, časové razítko a data transakce. Blockchain je decentralizovaná, distribuovaná a veřejná digitální účetní kniha, která je používána k zaznamenání transakcí napříč mnoha počítači, takže záznam nemůže být retroaktivně změněn bez úpravy všech následujících bloků a shody celé sítě. [1]

\obr{Vazba mezi blokem a hashem[3]}{fig:HashAndBlock}{1}{graphics/What-Is-Blockchain-Infographic.jpg}



\subsection{Historie}
V roce 1976 byla publikována práce "New Directions in Cryptography", která diskutovala koncept distribuované účetní knihy. S pokrokem v oblasti kryptografie byla publikována další práce nazvaná "How to Time-Stamp a Digital Document" od Stuarta Habera a Scotta Stornetty, která představila koncept časového razítka dat namísto média. Další důležitý koncept, který přispěl k vývoji blockchainu, byla "elektronická hotovost" nebo "digitální měna" založená na modelu navrženém Davidem Chaumem. Tento koncept následovaly například protokoly jako e-cash schémata, které detekovaly dvojité utracení. Adam Back v roce 1997 představil koncept "hashcash", který nabídl řešení pro kontrolu spamových e-mailů. To vedlo k vytvoření "b-money" Wei Daie založené na peer-to-peer síti. Považuje se, že Satoshi Nakamoto je vynálezcem technologie blockchain, když v roce 2008 publikoval dokument o bitcoinu "Bitcoin: A Peer-to-Peer Electronic Cash System". Jeho cílem bylo umožnit přímou online platbu z jednoho zdroje na druhý bez spoléhání se na třetí stranu. Jeho dokument navrhoval elektronický platební systém založený na kryptografii, který řešil problém dvojitého utracení, kde digitální měna nemůže být duplikována a nelze ji utratit více než jednou. [2]



\subsection{Principy}
Jedinečnost a krása blockchainu je v jejím distrubovaném přístupu a neměnnosti, která díky samotné nátuře znemožnuje podvody a to i bez nutnosti centralní entity.

\subsubsection{Decentralizace}
V kontextu blockchainu decentralizace představuje převedení dozoru a rozhodovací zodpovědnosti od centralizované entity (jednotlivec, společnost nebo skupina lidí) do distribuované sítě. Účelem těchto distribuovaných sítí je minimalizovat úroveň důvěry, kterou by měli účastníci vkládat jeden do druhého, a bránit v uplatňování autority nebo kontroly, která by mohla negativně ovlivnit efektivitu dané sítě.[5?]




\subsubsection{Immutability}
Klíčovou vlastností technologie blockchain je neměnnost. To v praxi znamená, že po vložení dat do sítě již není možno s těmito záznamy nijak manipulovat. Informace jsou uloženy v řetězci datových bloků, přičemž každý blok odkazuje na svého předchůdce pomocí hodnoty uložené v záhlaví ve tvaru alfanumerického řetězce, který je odvozen od předešlého bloku (typicky pomocí hash funkce) ,tak udržuje pospolitost bloků. Tato skutečnost znemožňuje provádění změn v datech předešlých bloků, protože by již nesouhlasily hodnoty odkazující na předešlé datové bloky. Němennost je jedním z důvodů proč technologie blockchain vzbuzuje velkou důvěru mezi uživateli. [4]

Neměnnost je vlastnost blockchainu, která zaručuje jeho věrohodnost a bezpečnost, ta je zaručena přímočarou logikou, která i tak čelí výzvám. Jednou z možných hrozeb je tzv. Útok 51, kdy útočník, nebo skupina útočníků ovládne více než 50 procent hashovací síly sítě a tak může blokovat nové transakce nebo je dokonce vracet, což vede k dvojtému utácení.[mozna 6] "Kvantové počítače mohou v budoucnosti také představovat problém, protože v blockchainu je využita asymetrická kryptografie, lze tak získat z podpisu transakce veřejný klíč a při dostatečném výpočetním výkonu uhádnou klíč soukromý. Tento problém se netýká ale výhradně blockchainu, nýbrž kryptografie a asymetrické kryptografie obecně." [mozna 6]

\subsubsection{Bloky}\label{subsubsec:Bloky}
Blok je kontejner, který uchovává sadu transakcí a další metadata. Skládá se ze dvou částí a to záhlaví a těla. Záhlaví se mohou v různých systémech lišit, zpravidla ale obsahují informace o samotném bloku a jejím autorovi vyjádřených těmito atributy : 

\begin{itemize}
	\item Hash předchozího bloku: Je kryptografický otisk využíván pro propojení bloku n+1 s blokem n pomocí hash kódu. Stručně řečeno, tato metoda slouží jako odkaz na hash předchozího (rodičovského) bloku v řetězci.[7]
	\item Nonce: zkratka z anglického "number used only once" (číslo použité pouze jednou), je náhodné číslo, které těžaři hledají s cílem vytěžit nový blok (více v sekci \ref{subsubsec:Mining})
	\item Časový otisk: anglicky "Timestamp" identifikuje čas, kdy byl blok vytěžen
	\item Náročnost: anglicky "Difficulty" je měřítko definující složitost prováděných operací v dané síti
	\item Merkle root hash: kořenový kryptografický otisk (nebo-li hash) všech transakcí v bloku. Více v sekci \ref{subsubsec:MerkeRootHash}
	\item Výška bloku: anglicky "Block height" inkrementovaná hodnota vyjadřující počet bloků vytěžených od počátku (anglicky "genesis")
\end{itemize}


\obr{Struktura bloků[7]}{fig:HashAndBlock}{1}{graphics/Structureofblocksinblockchain.png}


Tělo bloku obsahuje seznam všech transakcí. Každá transakce zahrnuje adresy odesílatele a příjemce, množství převáděné kryptoměny a digitální podpisy zapojených stran. Transakce v bloku jsou organizovány v určitém pořadí, kde prvním je speciální transakce označovaná jako "coinbase". Tato transakce představuje způsob, jakým jsou do systému zavedeny nové mince a jak je těžař odměněn za svou práci.[8]




\subsubsection{Transakce}

Transakci můžeme definovat jako smlouvu, dohodu, převod nebo výměnu aktiv mezi dvěma nebo více stranami. Obvykle se jedná o hotovost nebo majetek. [9]

Obdobně se transakce v blockchainu nejedná o nic jiného než přenos dat napříč sítí počítačů blockchainového systému. Síť počítačů blockchainu ukládá transakční data jako repliky v úložišti, kterému se často říká digitální účetní kniha. [9]

Uživatel inicializuje transakci v aplikaci a poté ji odesílá do blockchainu. Transakce uživatele je následně umístěna do "mempoolu" (dobře organizované fronty, ve které jsou transakce uschovávány a tříděny před přidáním do nově vytvořeného bloku) s dalšími čekajícími transakcemi a čeká na své "těžení". Jakmile dostatečné množství těžařů potvrdí transakci, je trvale uložena v bloku. Podle principu jsou transakce v blockchainu nevratné. Více v sekci \ref{subsubsec:Mining} [9]

Blok obsahuje sadu transakcí, má omezení jak časová, tak velikostní ve smyslu kolik transakcí může obsahovat. Po uplynutí časovače bloku nebo při dosažení maximálního počtu transakcí je blok přidán do řetězce a začíná iterace dalšího bloku.[9]





\subsubsection{Mining}\label{subsubsec:Mining}
Těžení je typické pro sítě implementující PoW (Proof-of-work) konsensus mezi které patří i nejznámější kryptoměna Bitcoin. Více o různých typech konsensů v kapitole \ref{subsec:Konsensus}

Těžba je proces, který Bitcoin a některé další kryptoměny využívají k vytváření nových mincí a ověřování nových transakcí. Zahrnuje velké množství decentralizované sítě počítačů po celém světě, které ověřují a zabezpečují blokchain. Jako odměnu za přispění svého výpočetního výkonu jsou těžaři v síti odměněni novými mincemi. Jedná se o ctnostný cyklus: těžaři udržují a zabezpečují blockchain, blockchain je odměňuje mincemi, mince poskytují podnět pro těžaře k udržování blockchainu.[10]





\subsection{Kryptografie}
Kryptografie představuje metodologii pro ochranu informací před neautorizovaným přístupem. V blockchain technologii je kryptografie využívána k zabezpečení transakcí prováděných mezi dvěma uzly. K tomuto slouží dva fundamentální koncepty: kryptografie a hashování (viz \ref{subsubsec:Hash}). Kryptografie se používá pro šifrování zpráv v P2P síti, hashování slouží k zabezpečení dat bloků a pro spojení bloků v blockchainu.

Kryptografie se primárně soustředí na zajištění bezpečnosti účastníků, transakcí a ochranu proti dvojitému utracení. Pomáhá zabezpečovat různé transakce v blockchainové síti. Zajišťuje, že pouze jednotlivci, pro které jsou data transakce určena, mohou získat, číst a zpracovávat transakci.[11]





\subsubsection{Symetrická kryptografie}\label{subsubsec:SymKrypt}

Tento typ šifrování používá stejný klíč pro oba procesy - šifrování i dešifrování. Symetrické šifrování je primárně používáno pro ochranu webových spojení a šifrování dat. Může být také označováno jako kryptografie s tajným klíčem. Jednou z výzev této metody je bezpečná výměna klíčů mezi odesílatelem a příjemcem. Data Encryption System (DES) je známým příkladem systému s symetrickým klíčem.[11]


\obr{Symetrická kryptografie [12]}{fig:SymEn}{1}{graphics/SymmetricEncrypt.jpg}


\subsubsection{Asymetrická kryptografie}

Metoda asymetrického šifrování používá různé klíče pro šifrování a dešifrování. K tomu využívá klíč veřejný a soukromí. Šifrování veřejným klíčem umožňuje sdílet informace s zcela neznámými stranami jako například ID Email. Soukromý klíč poté slouží k dešifrování zprávy a ověření digitálního podpisu (více v \ref{subsubsec:DigSig}). Matematický vztah mezi klíči spočívá v tom, že soukromý klíč nelze odvodit z veřejného klíče, ale veřejný klíč lze odvodit od klíče soukromého. Příklady asymetrické kryptografie jsou: ECC, DSS a další. [12]


\obr{Asymetrická kryptografie [12]}{fig:SymEn}{1}{graphics/AsymmetricEncrypt.jpg}





\subsubsection{Digitální podpis}\label{subsubsec:DigSig}


Digitální podpis je kryptografická technika používaná k ověření pravosti zprávy, transakce nebo dokumentu. 

V kontextu blockchainu se digitální podpis používá k prokázání vlastnictví určité transakce nebo dat jednotlivcem nebo entitou.

Digitální podpisy jsou zásadní součástí technologie blockchain. Zajišťují pravost a integritu transakcí potažmo dat v síti. Poskytují bezpečný a decentralizovaný způsob ověřování vlastnictví, eliminují potřebu prostředníků a zvyšují důvěru mezi stranami.[13]


\subsubsection{Digitální peněženka}\label{subsubsec:DigWallet}

Digitální nebo-li Blockchainová peněženka je speciální softwarové nebo hardwarové zařízení, které se používá k uchování informací o transakcích a osobních informací uživatele. Blockchainové peněženky neobsahují skutečnou měnu. Peněženky se používají k uchování soukromých klíčů a udržování bilance transakcí.

Peněženky jsou pouze komunikační nástroj pro provádění transakcí s ostatními uživateli. Skutečná data nebo měna jsou uložena v blocích v blockchainu.[11]







\subsubsection{Hashovací funkce}\label{subsubsec:Hash}

Hash je matematická funkce, která převádí vstup s libovolnou délkou na šifrovaný výstup s pevnou délkou. Tudíž, bez ohledu na původní množství dat nebo velikost souboru, bude jeho jedinečný hash vždy stejné velikosti. Co je však důležité, hash nelze použít k \it{reverse-engineeringu}\textsuperscript{1} vstupu z hashovaného výstupu, protože hashovací funkce jsou „jednosměrné“. To znamená, že pokud použijete takovou funkci na stejná data, její hash bude identický, takže pokud již znáte jeho hash, můžete ověřit, že data jsou stejná (tedy nezměněná).

Hashovací funkce jsou základem pro vytváření odtisků dat (tzv. checksumů) a zajišťují integritu a autenticitu dat v rámci blockchainu. Dále, v kontextu blockchainových technologií, hashování je zásadní pro vytváření nových bloků. Pokud je jakýkoliv prvek bloku změněn, změní se celkový hash bloku, což efektivně znamená, že jakákoliv neautorizovaná změna je okamžitě viditelná. Takové vlastnosti hashovacích funkcí přidávají vrstvu bezpečnosti a transparentnosti do technologie blockchain.[14]






\subsubsection{Merke root hash}\label{subsubsec:MerkeRootHash}

Merkle Tree je základním prvkem technologie blockchain. Jedná se o matematickou datovou strukturu složenou z hashů jednotlivých datových bloků, které slouží jako shrnutí všech transakcí v bloku. Tato struktura umožňuje efektivní a bezpečné ověřování obsahu v rozsáhlém množství dat a tak pomáhá při ověřování konzistence a samotného obsahu dat. Jak Bitcoin, tak Ethereum používají strukturu Merkle Tree. Merkle Tree je také známý jako Hashovací strom. 

Koncept Stromu Merkle je pojmenován po Ralphu Merkleovi, který tento nápad patentoval v roce 1979.[15]


\obr{Zjednodušená forma Merkle Tree [15]}{fig:MerkleTree}{0.6}{graphics/blockchain-merkle-tree.png}




\subsubsection{Výhody a limitace}

Výhody:
\begin{itemize}
	\item Přístupnost: Jednou z hlavních výhod technologie blockchain je její otevřenost. To znamená, že se kdokoli může stát účastníkem a přispět. Pro připojení k distribuované síti není nutná žádná smlouva nebo specifické povolení.
	
	\item Decentralizace: Technologie blockchain je považována za osvobozenou od cenzury, jelikož není ovládána žádnou jednotlivou entitou. Spíše funguje na principu důvěryhodných uzlů pro ověřování a konsenzuálních protokolů, které schvalují transakce pomocí chytrých kontraktů.

	\item Bezpečnost: Technologie blockchain dosahuje vysoké úrovně bezpečnosti díky využití hashovacích technik. Ty umožňují ukládání každé jednotlivé transakce do bloků, které jsou vzájemně propojeny. To zajišťuje robustní ochranu dat. Jednou z takových technik, kterou blockchain užívá pro ukládání transakcí, je SHA 256.[16]

	
\end{itemize}


Nevýhody:

\begin{itemize}
	
	\item Škálovatelnost: Jedna z největších nevýhod technologie mnoha technologií blockchain technologií je pevná velikost bloku, která je 1 MB. To omezuje počet transakcí, které může blok obsahovat a výrazně snižuje škálovatelnost a univerzálnost použití blockchainu.
	
	
	\item Časová náročnost: Proces přidání nového bloku do blockchainu vyžaduje, aby těžaři opakovaně vypočítávali hodnoty nonce (vysvětleno v sekci \ref{subsubsec:Bloky}). Tento proces je nesmírně časově náročný a je třeba ho optimalizovat, aby blockchain byl efektivně využitelný na průmyslové úrovni. Tento problém se týká výhradně řešeních využívající Proof-of-work konsensus.
	
	
	\item Regulační výzvy: Regulační výzvy se týkají technologie blockchain, která implementuje arbitrární monetární politky a tak čelí určitým výzvám, zejména ze strany finančních institucí. Pro širší uplatnění blockchainu je nutné řešit další technologické aspekty a navrhnout vhodné regulace.[16]
	
	
\end{itemize}



\subsection{Konsensus}\label{subsec:Konsensus}


\n{3}{Proof of Work}
text
\n{3}{Proof of Stake}
text
\n{3}{Ostatní algoritmy konsensu}
text



\n{1}{Decentralizované aplikace}
text
\n{2}{Web2 vs Web3}
text
\n{2}{Smart kontrakty}
text
\n{2}{Distribuované úložiště}
text
\n{2}{Výhodné a nevýhodné vlastnosti}
text
\n{2}{Sítě typu peer-to-peer}
text
\n{2}{Persistence}





\n{1}{Ethereum}
text
\n{2}{Ether}
text
\n{2}{Ucty}
text
\n{2}{Transakce}
text
\n{2}{Bloky}
text
\n{2}{Ethereum Virtual Machine}
text
\n{2}{Gas}
text
\n{2}{Uzly a klienti}
text
\n{2}{Konsensus}
text
\n{2}{Ethereum stack}
text
\n{3}{Smart kontrakty}
text
\n{3}{Development networks}
text




% \n{1}{Další nadpis}
% Tato sekce obsahuje ukázku vložení obrázku (Obr. \ref{fig:logo}).

% Obrázek lze vkládat pomocí následujícího zjednodušeného stylu, nebo klasickým LaTex způsobem
% Pozor! Obrázek nesmí obsahovat alfa kanál (průhlednost). Jde to proti standardu PDF/A.
\obr{Popisek obrázku}{fig:logo}{0.5}{graphics/logo/fai_logo_cz.png}

\n{2}{Podnadpis}
Tato sekce obsahuje ukázku vložení tabulky (Tab. \ref{tab:priklad}).

% Tabulku lze vkládat pomocí následujícího zjednodušeného stylu, nebo klasickým LaTex způsobem
\tab{Popisek tabulky}{tab:priklad}{0.65}{|l|c|c|c|c|c|r|}{
  \hline
   & 1 & 2 & 3 & 4 & 5 & Cena [Kč] \\ \hline
  \emph{F} & (jedna) & (dva) & (tři) & (čtyři) & (pět) & 300 \\ \hline
}

\n{3}{Podpodnadpis}

\n{3}{Podpodnadpis}
Citace knihy. \cite{chmel}


% ============================================================================ %

% Pokud Vaše práce neobsahuje analytickou část, stačí odstranit či zakomentovat nasledujících pár rádků
% \cast{Analytická část}

% \n{1}{Nadpis}

% \n{2}{Podnadpis}


% ============================================================================ %
\cast{Projektová část}

\n{1}{Aplikace pro ověřování dokumentů}
Text
\n{2}{Řešený problém}
text
\n{2}{Navrhované řešení}
text
\n{2}{Terminologie}
text
\n{3}{Dokument}
text
\n{3}{Atestace}
text


\n{1}{Nástroje a technologie pro vývoj}
text

\n{2}{Ethereum Attestation Service}
text
\n{3}{Atestace}
text
\n{3}{Schémata}
text
\n{3}{Onchain vs Offchain}
text
\n{3}{Privacy}
text
\n{3}{Resolver Contracts}
text
\n{3}{Developer Tools}
text
\n{3}{EAS SDK}
text


\n{2}{MetaMask (nebo jiná digitální peněženka)}
text
\n{3}{Integrace do aplikace}
text
\n{3}{API}



\n{2}{Inter Planetary File System}
text
\n{3}{IPFS vs HTTP}
text
\n{3}{Životní cyklus dat}
text
\n{3}{Hashovani}
text
\n{3}{Imutabilita}
text
\n{3}{Persistence}
text
\n{3}{Privacy a sifrovani}
text
\n{3}{Uzly}
text
\n{3}{Pinata}
text
\n{3}{Filecoin}
text
\n{3}{Web3.storage}
text
\n{3}{Fleek storage}
text

\n{2}{Front-end}
text
\n{3}{Zduvodneni volby}
text

\n{2}{Back-end}
text
\n{3}{Zduvodneni volby}




\n{1}{Implementace a postup vývoje}
text
\n{2}{Datový model}
\n{3}{Jake data sbirat}
\n{3}{Kam ukladat}
\n{3}{Jak ukladat}

\n{2}{Aktoři}
\n{3}{Odesilatel atestace}
\n{3}{Prijmce atestace}
\n{3}{Verifier}
kdo je to uzivatel aplikace

\n{2}{Use cases}
jake jsou use-cases

\n{2}{Konfigurace aplikace}
text
\n{3}{SDK EAS}
text
\n{3}{SDK MetaMask/neco jineho nevim zatim}
text
\n{3}{SDK IPFS}
text





\n{2}{Webová prezentace}
text
\n{2}{Video prezentace}



% ============================================================================ %
\nn{Závěr}
Text závěru


% ============================================================================ %


% PLAN PRACE NA BP

% soucasny plan:
% PRAKTICKA CAST
% v 1. kroku analyza technologii a nastroju {Ethereum(EVM, Smart Contracts, Security,Network,....), BC(consenus, co je to blok, transakce,...),Kryptologie(asymetricke sifrovani, hash funkce,.....),IPFS(jak ulozit dokument), MetaMask(intro do SDK, jak integrovat do aplikace), Ethereum Attestation Service(co to atestace, schemata, resolver contracts,....), dApps(co to je, )}
% v 2.kroku analyza dat, aktoru (jake data sbirat, kam ukladat, kdo je to uzivatel aplikace, jake jsou use-cases)
% v 3. kroku bude pak tech design (kde bych chtel rozebrat podobrne flow aplikace, popsat jak do sebe jednotlive SDK zapadaji, jak spolu komunikuji, poukazat na mozne vyzvy v implementaci, navrh FE).
% ve 4. kroku bych chtel udelat MVP(zakladni flow, jednoduchy FE)
% krok 5. zacit psat tu Praktickou cast ( popsat jak jsem postupoval pri tvorbe aplikace)
% TED ZACINA TEORETICKA CAST
% krok 6. teoreticka cast, napsat o blockchainu, ethereu, EAS, IFPS, krypotologii v BC, oduvodnit volbu technologii atd. % Hlavni text prace

\OdsazovaniOdstavcuStop


% ============================================================================ %
\seznamlit{
  % Na toto místo vložit veškeré citované bibliografické položky.
  	%\bibitem{chmel}\it{Český chmel: atlas odrůd} [online]. [cit. %2004-10-15]. Dostupný z~WWW: \url{http://www.beer.cz/humulus/}.
  
   	\bibitem{blockchainobecne}\it{blockchain obecne} [online]. [cit. 2004-10-15]. Dostupný z~WWW: \url{https://www.synopsys.com/glossary/what-is-blockchain.html}.

	\bibitem{blockchainhistory}\it{blockchain history} [online]. [cit. 2004-10-15]. Dostupný z~WWW: \url{https://d1wqtxts1xzle7.cloudfront.net/60715489/Understanding\_Blockchain\_Technology20190926-26770-147v9qi-libre.pdf?1569545201=\&response-content-disposition=inline/%3B+filename%3DUnderstanding_Blockchain_Technology.pdf&Expires=1710887337&Signature=CC7eg5Ia-e1qIVZ~mC79PtM9TYCLOoKQ-RvV9OlmggMZjw1pPbKxsZl-bKBDt1OPnHBe64BqPu4En-1XKfQ8lr-r-yUjA7AKwmpxHN6vt3X6pP7pNlGA62LkoRnddD0~eTSwAfwIXs0cqymOI6yLR6r8R9983teaXYZjnGz1BupFndGPzY6Ff2HxLXb91KKM1RqSCMm3GRpIyFTwOiA13aFNxfxxd6v8H14ZSE0X3wkLZ0zkzoZMWdeinrWD1xsg~DZ9XT3mZ4O8SvwrNO3JxmPINMEiTppBfVjc3V1CLLz6wC50SlVk6vaY97pEgtMdsmeWMWT1gKErFuuz7STy-g__&Key-Pair-Id=APKAJLOHF5GGSLRBV4ZA}.
	}

	\bibitem{blockAndHash}\it{Vazba bloku a hashe} [online]. [cit.2024-03-19]. Dostupný z~WWW:\url{https://money.com/what-is-blockchain/}
	
	\bibitem{Imutabilita blockchainu}\it{Imutabilita blockchain} [online]. [cit.2024-03-19]. Dostupný z~WWW:\url{https://www.researchgate.net/profile/Moritz-Boehmecke-Schwafert/publication/322322101_The_immutability_concept_of_blockchains_and_benefits_of_early_standardization/links/6420b404315dfb4cceaee61f/The-immutability-concept-of-blockchains-and-benefits-of-early-standardization.pdf}

	\bibitem{Decentralizace blockchainu}\it{Decentralizace blockchainu} [online]. [cit.2024-03-19]. Dostupný z~WWW:\url{https://www.blockchain-council.org/blockchain/what-is-decentralization-in-blockchain/}
	
	\bibitem{hrozby immutability}\it{hrozby immutability} [online]. [cit.2024-03-19]. Dostupný z~WWW:\url{https://blog.cfte.education/immutable-ledger-in-blockchain/}
	
	\bibitem{struktura bloku}\it{struktura bloku} [online]. [cit.2024-03-19]. Dostupný z~WWW:\url{https://www.geeksforgeeks.org/blockchain-structure/}
	
	\bibitem{telo bloku}\it{telo bloku} [online]. [cit.2024-03-19]. Dostupný z~WWW:\url{https://www.theblock.co/learn/245697/what-are-blocks-in-a-blockchain}
	
	\bibitem{Transakce}\it{Transakce} [online]. [cit.2024-03-19]. Dostupný z~WWW:\url{https://www.scalingparrots.com/en/what-is-a-blockchain-transaction/
	}
	
	\bibitem{Mining}\it{Mining} [online]. [cit.2024-03-19]. Dostupný z~WWW:\url{https://www.coinbase.com/en-gb/learn/crypto-basics/what-is-mining}

	\bibitem{Kryptografie v Blockchain}\it{Kryptografie v Blockchain} [online]. [cit.2024-03-19]. Dostupný z~WWW:\url{https://www.geeksforgeeks.org/cryptography-in-blockchain/}
	
	\bibitem{Sifrovani}\it{Sifrovani} [online]. [cit.2024-03-19]. Dostupný z~WWW:\url{https://www.analyticsvidhya.com/blog/2022/09/concept-of-cryptography-in-blockchain/}
	
	
	\bibitem{Digitální podpis}\it{Digitální podpis} [online]. [cit.2024-03-19]. Dostupný z~WWW:\url{https://mudrex.com/learn/digital-signature-blockchain-importance/}
	
	
	\bibitem{Hash funkce}\it{Hash funkce} [online]. [cit.2024-03-19]. Dostupný z~WWW:\url{https://www.investopedia.com/terms/h/hash.asp}
	
	
	\bibitem{Merkle Tree}\it{Merkle Tree} [online]. [cit.2024-03-19]. Dostupný z~WWW:\url{https://www.javatpoint.com/blockchain-merkle-tree}
	
	
	
	\bibitem{Výhody a nevýhody BC}\it{Výhody a nevýhody BC} [online]. [cit.2024-03-19]. Dostupný z~WWW:\url{https://www.geeksforgeeks.org/advantages-and-disadvantages-of-blockchain/}
	
	
	
	
	
}

% Pro generování literatury lze alternativně použít i příkaz "\seznamlitbib", 
% který se postará o plnohodnotné vkládání referencí pomocí "bibliography". 
% V takovém případě se využívají bibliografické údaje uložené v souboru 
% tex-literatura.bib. Ty se automaticky upravuji dle zvolené citační normy 
% (v šabloně je nastavena korektní česká norma).
%\seznamlitbib


% ============================================================================ %
\input{tex/zkratky.tex} % Seznam zkratek


% ============================================================================ %
\seznamobr  % Seznam je generován automaticky


% ============================================================================ %
\seznamtab  % Seznam je generován automaticky


% ============================================================================ %
\input{tex/prilohy.tex} % Prilohy


% ============================================================================ %

\end{document}

% ============================================================================ %
